\section{Functional Requirements}
\label{chap:requirements}
To organize and prioritize features in Hover, we have created three sets of functional requirements (FR) based on their importance; \textit{High}, \textit{medium}, and \textit{low}. This chapter will present and briefly explain the idea behind the categories.

\subsection{High Priority FRs}
High-priority Functional Requirements are needed for the game to function as intended. Consequently, the requirements listed below are the ones that we categorize as a must-have in the application before testing it on real users. Thus, all eight high-priority FRs are implemented in Hover.
\begin{itemize}
    \item \textbf{FR1:} The user should be able to create an account and log in.
    \item \textbf{FR2:} The user should be able to view and explore hover zones (i.e., geofences).
    \item \textbf{FR3:} The user should be able to track and publish activities.
    \item \textbf{FR4:} The user should be able to challenge friends.
    \item \textbf{FR5:} The user should be able to view friends’ activities, earned achievements, and completed challenges.
    \item \textbf{FR6:} The user should be able the compare their score with friends on a leaderboard.
    \item \textbf{FR7:} The user should be able to earn achievements for completing predefined tasks.
    \item \textbf{FR8:} The user should be able to increase an activity score by socializing.
\end{itemize}

\subsection{Medium Priority FRs}
Medium priority FRs are requirements that should be implemented if we have time. They improve the game experience but are not crucial for the game to function as intended. Three of the medium priority FRs, FR9, FR10, and FR11, have been implemented.

\begin{itemize}
    \item \textbf{FR9:} The user should get notifications of important events.
    \item \textbf{FR10:} The user should be able to create a custom profile and visit their friend’s profile.
    \item \textbf{FR11:} The user should be able to interact with friends’ activities, earned achievements, and completed challenges.
    \item \textbf{FR12:} The user should be able to customize their list of friends.
\end{itemize}

\subsection{Low Priority FRs}
Low priority functional requirements are nice to have features. They could improve the game, but not in a dramatic way. In this project, the requirements categorized with low priority have not been implemented. However, they serve as good suggestions for further work, as we will describe later in Chapter \ref{chap:further_work}. 

\begin{itemize}
    \item \textbf{FR13:} The user should be able to earn real-world awards.
    \item \textbf{FR14:} The user should be able to evolve a character based on their score.
    \item \textbf{FR15:} The user should be able to log in using external services, Facebook, Google, Contacts, etc.
\end{itemize}

\subsection{Summary}
In this chapter, the functional requirements have been listed and categorized. The requirements are used when describing the game in Chapter \ref{chap:game_description} and explaining game design elements in Chapter \ref{chap:game_design}. Chapter \ref{chap:game_development} also includes a more detailed explanation of how the requirements have been used during the development process.